\section{Response Availability and Gilbert--Lynch}
\label{app:cap-response}

Complete CAP (Theorem~\ref{thm:complete-cap}) asks that every permitted output
remain available during a partition; Gilbert and Lynch ask only that every
request receive some response~\cite{gilbert2002cap}.  This appendix defines the
second notion, proves the corresponding form of Complete CAP, and derives the
Gilbert--Lynch theorem.  Throughout, fix an ordered relational specification
$\Spec$ and the partitioned histories of Section~\ref{sec:complete-cap}.

\subsection{Obligations, availability, and locality}

\begin{definition}[Response obligation]
\label{def:response-obligation}
A \emph{response obligation} assigns to each participant $p$ a relation
$Q_p\subseteq S$ that is total on $I$ and upward closed within fibers: if
$(i,o)\in Q_p$, $(i,u)\in S$, and $o\OutputOrd u$, then $(i,u)\in Q_p$.  An
output in $Q_p(i)$ \emph{answers} at $i$ for $p$.  The obligation is
\emph{uniform} when every $Q_p$ is the same relation $Q$.
\end{definition}

The obligation is part of the declared specification.  Response availability
is only as strong as $Q$: if an output permitted at every input answers, every
specification is available.

\begin{definition}[Response Availability]
\label{def:response-availability}
A system has \emph{Response Availability} for $(Q_p)_{p\in P}$ when, at every
consistent cut of every partitioned history it produces, each participant $p$
with current input $i$ satisfies two conditions, each by a continuation that
the system produces, that is partitioned by the same components, and that adds
only events of $p$'s component.
\begin{enumerate}
  \item \emph{Answering.}  In a continuation that adds no input event at $p$,
        $p$ emits some output in $Q_p(i)$.
  \item \emph{Input acceptance.}  For every $j\InputGe i$, in some
        continuation $p$ records $\mathsf{input}(j)$.
\end{enumerate}
\end{definition}

Input acceptance says that $p$ never seals an \indesc of its current input.
It excludes, for example, the election under the equality order
(Figure~\ref{fig:nonmonotone-frame}), and it matches Gilbert and Lynch's
requirement that nodes accept every request.

\begin{definition}[Response Partition Consistency]
\label{def:response-partition-consistency}
A consistent cut has \emph{Response Partition Consistency} at an input $j$
that \inrel every participant's current input when some $u\in\fiber{j}$
\outrel every output emitted in the cut.  A system has it when every
consistent cut of every partitioned history it produces has it at every such
$j$.
\end{definition}

Confluence at $j$ implies Response Partition Consistency at $j$
(Corollary~\ref{cor:reconciliation-any-basis}), and for a canonical system
(Definition~\ref{def:canonical-system}) the two coincide
(Lemma~\ref{lem:reconciliation}).  With other commitments the converse can
fail: participants may prune their states so that they cannot meet while their
emitted outputs still have a common \outnoun.

A system is \emph{replicated} when its obligation is uniform and every
participant runs the same strategy, which depends only on the participant's
local history and not on its identity or the number of participants.  The
environment supplies any input to any participant and resolves the
strategies' nondeterministic choices.

\begin{lemma}[Locality]
\label{lem:locality}
Let local runs that the strategies of a system allow, one per participant,
each from that participant's initial input and none containing a receive, be
placed in separate components, with any further participants idle.  The result
is a partitioned history that the system produces.  If the system is
replicated, a local run that the strategy allows for one participant is
allowed for any participant with the same initial input, in any participant
set.
\end{lemma}

\begin{proof}
Whether a history is produced is checked event by event against each
participant's own local prefix, and there are no receives to match.
Replication makes that check independent of the participant.
\end{proof}

\subsection{The response form of Complete CAP}

\begin{definition}[Response interface]
\label{def:response-interface}
For a uniform obligation $Q$, a \emph{response interface} is a relation
$T\subseteq Q$ that is total on $I$ and \emph{self-continuing}: if
$(i,o)\in T$ and $i\InputOrd j$, then some $u\in T(j)$ \outrel $o$.  The
specification is \emph{consistent on $T$} when, for every finite set of rows
$(i_1,o_1),\ldots,(i_k,o_k)$ of $T$ and every common \indesc $j$ of
$i_1,\ldots,i_k$, some $u\in\fiber{j}$ \outrel every $o_r$.
\end{definition}

The reconciling output may lie outside $T$.  With $T=S$, consistency on $T$ is
the common-\indesc condition of Proposition~\ref{prop:common-expansion-form},
and by Lemma~\ref{lem:reconciliation} consistency on $T$ is confluence of the
canonical runs whose emitted rows lie in $T$.

\begin{theorem}[Complete CAP, response form]
\label{thm:complete-cap-response}
For an ordered relational specification $\Spec$ and a uniform obligation $Q$,
the following are equivalent:
\begin{enumerate}
  \item some replicated system has Response Availability for $Q$ and
        Response Partition Consistency;
  \item $\Spec$ is consistent on some response interface $T\subseteq Q$.
\end{enumerate}
\end{theorem}

\begin{proof}
(2$\Rightarrow$1.)  Take the canonical basis and the strategy that ignores
messages, records every input, and emits by the canonical event some $u\in
T(i)$ at its current input $i$ that \outrel its last emitted output, or any
$u\in T(i)$ if it has emitted nothing.  Self-continuation supplies $u$.  The
emitted outputs form a chain $o_1\OutputOrd\cdots\OutputOrd o_m$, so by
Lemma~\ref{lem:run-residual} the table is $S[\mathrm{in}\InputGe i,\
\mathrm{out}\OutputGe o_m]$ (or $S[\mathrm{in}\InputGe i]$), which contains
$(i,u)$: this is Answering.  For $k\InputGe i$, self-continuation applied to
the row of $o_m$ (or totality of $S$) leaves a row at $k$: this is Input
acceptance.  At a consistent cut, the emitted rows are finitely many rows of
$T$ whose inputs every common \indesc $j$ of the current inputs \inrel, and
consistency on $T$ gives Response Partition Consistency.

(1$\Rightarrow$2.)  Call an emission \emph{isolated} when its participant has
received no message before it, and let $T$ be the set of rows $(i,o)$ with
$o\in Q(i)$ that some participant emits in an isolated emission, at current
input $i$, in some partitioned history of the system.  Then $T\subseteq Q$.
By Lemma~\ref{lem:locality}, each witness can be replayed by a participant
alone in its component, with an idle participant forming a second component.
\emph{Totality:} at the empty cut, Answering gives an isolated emission at the
initial input.  \emph{Self-continuation:} given $(i,o)\in T$ and $j\InputGe
i$, replay the witness, record $\mathsf{input}(j)$ by Input acceptance, and
answer by Answering with some $u\in Q(j)$.  By
Lemmas~\ref{lem:canonical-equivalent} and~\ref{lem:run-residual} the table
before that emission is contained in $S[\mathrm{in}\InputGe j,\
\mathrm{out}\OutputGe o]$, so $o\OutputOrd u$ and $(j,u)\in T$.
\emph{Consistency on $T$:} given rows $(i_r,o_r)$ of $T$ with common \indesc
$j$, replay each witness on its own participant, truncated after the
emission, and add a participant with initial input $j$.  By
Lemma~\ref{lem:locality} this is a partitioned history of the system, and
Response Partition Consistency at its final cut gives $u\in\fiber{j}$
\outrelpart every $o_r$.  The empty family is covered by totality of $S$.
\end{proof}

Without replication, the same argument gives one response interface $T_p\subseteq
Q_p$ per participant, with consistency required only for families that use each
$T_p$ at most once.

\subsection{Relationship to the interface form}

Lifted Monotonicity implies the response form for every uniform obligation,
with $T=Q$.  If $(i,o)\in Q$ and $i\InputOrd j$, Relational Monotonicity gives
$u\in\fiber{j}$ with $o\OutputOrd u$, totality gives some $q\in Q(j)$,
Directed Outputs gives $w\in\fiber{j}$ \outrelpart $u$ and $q$, and upward
closure puts $w$ in $Q(j)$; the common-\indesc condition gives consistency on
$Q$.

The converse fails.  Let $I=\{0\InputOrd 1\}$, $\bot\OutputOrd a$ and
$\bot\OutputOrd b$ with $a,b$ incomparable, $\fiber{0}=\{\bot,a\}$,
$\fiber{1}=\{b\}$, and $Q=S$.  Lifted Monotonicity fails at $(0,a)$, but $\Spec$
is consistent on the response interface $\{(0,\bot),(1,b)\}$.

For a functional specification the forms coincide: a total $T\subseteq Q\subseteq
S$ equals $S$, and both self-continuation and consistency on $S$ are
monotonicity.

A total $T\subseteq Q$ that satisfies Lifted Monotonicity as a specification
is a response interface on which $\Spec$ is consistent, but the condition is
weaker.  Let $i_1,i_2$ be incomparable inputs below $j$, and $k_1,k_2$
incomparable inputs above $j$ with no common \indesc.  Let
$\fiber{i_1}=\{a\}$, $\fiber{i_2}=\{b\}$, $\fiber{j}=\{c_1,c_2,d_1,d_2\}$,
$\fiber{k_1}=\{e_1\}$, $\fiber{k_2}=\{e_2\}$, $Q=S$, with $a\OutputOrd c_1$,
$b\OutputOrd c_2$, $d_1,d_2$ above $c_1$ and $c_2$, $e_1$ above $c_1,c_2,d_1$
but not $d_2$, and $e_2$ above $c_1,c_2,d_2$ but not $d_1$.  Then
$T=\{a,b,c_1,c_2,e_1,e_2\}$ is a response interface on which $\Spec$ is
consistent ($d_1$ reconciles at $j$, $e_1$ or $e_2$ at $k_1$ or $k_2$).  A
sub-specification with Lifted Monotonicity must contain $a$ and $b$, hence
$d_1$ or $d_2$ at $j$, and either has no \outnoun at $k_2$ or $k_1$.  The
example also shows that reconciling outputs need not lie in $T$: a participant
that emits $d_1$ or $d_2$ seals $k_2$ or $k_1$.

\subsection{Recovering Gilbert--Lynch}

\jmh{Consider option (b): promote the footnote's per-replica (non-uniform obligation) version of the response-form theorem to a theorem with proof, and derive Gilbert--Lynch from it, so that GL is a special case of a CAP theorem rather than a corollary of Lemma forced-answer.}

\begin{lemma}[Forced answer]
\label{lem:forced-answer}
Let $i\InputOrd j$, and suppose no output in $Q_q(i)$ has a \outnoun in
$\fiber{j}$.  In a system with Response Availability, consider a partitioned
history in which $q$ is alone in its component, with a cut at which $q$'s
current input is $i$ and $j$ \inrel every participant's current input.  Then
a continuation that adds only events of $q$, and no input event, reaches a cut
at which the emitted outputs have no common \outnoun in $\fiber{j}$.
\end{lemma}

\begin{proof}
Answering gives such a continuation in which $q$ emits some $o\in Q_q(i)$.  No
current input changes, so $j$ still \inrel all of them, and no output of
$\fiber{j}$ \outrel $o$.
\end{proof}

\paragraph{The register.}
An input is a finite operation history: a set of read and write invocations at
replicas with a strict partial order of real-time precedence.  One history
\inrel another when it contains all of its operations and precedence edges.
An output is a response map assigning values to reads and acknowledgments to
writes, ordered by extension.  The fiber $\fiber{h}$ contains the maps
explained by some linearization of $h$, a total order of its operations
extending precedence~\cite{herlihy1990linearizability} in which each answered
read returns the last preceding write, or the initial value $v_0$.  The
obligation $Q_p(h)$ consists of the maps in $\fiber{h}$ that answer every read
and acknowledge every write invoked at $p$.  Every finite history has a
linearization, so $Q_p$ is total, and it is upward closed within fibers.
Precedence is input because the environment sees a response before invoking
the next operation (Section~\ref{sec:election-story}).

Real time is the order of events in an execution: an \emph{interleaving}
$\lambda$ of a distributed history is a linear order of its events consistent
with happens-before, and no strategy observes it.  The \emph{invocation} of an
operation is the first input event at its replica whose input contains it,
and its \emph{response} is the first output event at that replica whose map
answers it.  For a prefix $C$ of $\lambda$ (a consistent cut), the
\emph{realized history} $h_C$ contains the operations invoked in $C$, with an
edge from $a$ to $b$ when $a$'s response precedes $b$'s invocation in
$\lambda$.  If $C'\subseteq C$ is also a prefix, $h_C$ \inrel $h_{C'}$.  The
realized history need not be a common \indesc of the participants' current
inputs: an input event may introduce two operations with an edge between them
before either responds.

\begin{corollary}[Gilbert--Lynch]
\label{cor:gilbert-lynch}
No system for the linearizable register has Response Availability and
Response Partition Consistency at $h_C$ for every prefix $C$ of every
interleaving of every partitioned history at which $h_C$ is a common \indesc
of the participants' current inputs.
\end{corollary}

\begin{proof}
Put replicas $p$ and $q$ in separate components from the start, and invoke a
write of \dog{} at $p$.  Answering gives a continuation in which $p$
acknowledges it.  Then invoke a read $r_q$ at $q$, after the acknowledgment in
$\lambda$.  No message crosses the partition, so $q$'s current input is
$i_q=\{r_q\}$, and $j=h_C$ contains the write, $r_q$, and an edge from the
write to $r_q$; it is a common \indesc of both current inputs.  Every map in
$Q_q(i_q)$ answers $r_q$ with $v_0$, and every linearization of $j$ places the
write first, so no output in $Q_q(i_q)$ has a \outnoun in $\fiber{j}$.
Lemma~\ref{lem:forced-answer} gives a continuation adding no input event, so
the realized history stays $j$, whose emitted outputs have no common
\outnoun in $\fiber{j}$.
\end{proof}

The register's obligation differs by replica, so
Theorem~\ref{thm:complete-cap-response} does not apply as stated.  The
correspondence with Gilbert and Lynch's asynchronous model has three parts.
\begin{enumerate}
  \item \emph{Algorithms are strategies.}  A replica's algorithm is a
        possibly nondeterministic function of its local history; its messages
        are messages, an invocation is input, and a response is an output
        event; we take the canonical basis.  The partition holds from the
        start, as in their witness executions, so a replica alone in its
        component cannot distinguish the execution from one in which the
        other component receives no invocations.  Atomicity there places each
        response in the fiber at the replica's view, so restricting responses
        to such outputs loses no algorithm that is atomic in every execution.
        This is where their indistinguishability argument lives.
  \item \emph{Availability} is Answering: every request to a non-failed node
        is answered, and messages across the partition are lost.  The
        corollary uses only Answering.
  \item \emph{Atomicity} is Response Partition Consistency at the realized
        history.
\end{enumerate}
